% Focal closure of antecedents: sources from Article XI, edition of 21 September 2026.
% Ranges, complete copies, and hashes: gestion/mapa_restitucion_memoria.json.
\section{Antecedents of operator reconstruction and memory}
\label{md:antmem:capitulo}

These antecedents begin after the joint APP--TRIT--TPK production
of the enriched state and the discrete structure of the
continuum. We work with the dodecaphase record already generated and with
declared realizations of its readout operators. The subsequent linear carrier
and the candidates in a reconstruction do not replace the domain of
generated histories. Nor is a subsequent reversible realization
identified with nonadic holonomy \(\Gamma_9\) merely through its notation.

In the chart and orientation of the record,
\begin{equation}
 \begin{gathered}
 K=(234,543,140,729,659,824,621,58,914,794,146,601),\\
 \mathbf1^{*}K=6263,\qquad
 (Sx)_i=x_{i+1},\quad i\pmod {12}.
 \end{gathered}
 \label{md:antmem:registro}
\end{equation}
The record is used after its generation, not as a retrospective input
for selecting an output. Here the operator \(S\) is
the marked positional shift. Its domains will remain
distinct from the canonical planes used later.

\subsection{The full orbit as a reference for its carrier}
\label{md:antmem:k:orbita}

In \(V=\mathbb R^{12}\), with its canonical Euclidean inner product, consider
\[
 O_K=(K,SK,\ldots,S^{11}K):\mathbb R^{12}\longrightarrow V.
\]
The shift \(S\) acts on positions in the record. Periodicity alone
does not identify it with the full TPK evolution or with
a subsequent exceptional symmetry.

\begin{theorem}[Complete cyclic reference and exact Gram matrix]
\label{md:antmem:k:marco-ciclico}
The matrix \(O_K\) is invertible. Its Gram matrix has eigenvalues
\[
\begin{array}{c|c}
 \text{value}&\text{multiplicity}\\\hline
 39225169=6263^2&1\\
 1248025-345420\sqrt3&2\\
 589609&2\\
 1407529&2\\
 1053157&2\\
 1248025+345420\sqrt3&2\\
 697225&1
\end{array}
\]
and satisfies
\begin{equation}
 589609\|c\|^2\leq\|O_Kc\|^2\leq39225169\|c\|^2,
 \quad
 \|O_K^{-1}\|=\frac1{\sqrt{589609}},\quad
 \operatorname{cond}(O_K)=\frac{6263}{\sqrt{589609}}.
 \label{md:antmem:k:cotas-orbita}
\end{equation}
The orbit of the centred record has rank eleven.
\end{theorem}
\begin{proof}
The Gram entries are the autocorrelations
\(\langle K,S^jK\rangle\). In order \(j=0,\ldots,11\), they are
\begin{align*}
 (&4251257,2999323,3163385,3287920,3216553,3344743,\\
  &2950064,3344743,3216553,3287920,3163385,2999323).
\end{align*}
Since \(S\) is orthogonal, the Gram matrix is circulant. Substituting the twelve
characters \(\exp(2\pi i m/12)\) into that row gives the preceding
real table. All values are positive. The smaller of the two values
involving \(\sqrt3\) exceeds \(589609\), since
\(658416>345420\sqrt3\), an inequality verified by squaring
its two positive sides. The other values in the table are
at least \(589609\). The uniform mode has value
\((\sum K_i)^2=6263^2\); the triangle inequality for the Fourier
transform, using \(K_i\geq0\), bounds all the others by that same
value. The spectral theorem gives \eqref{md:antmem:k:cotas-orbita}. Centring
the record removes only the uniform mode; the remaining eleven retain
the nonzero values in the table.
\end{proof}

The determinant, with the column order and orientation of \(S\)
fixed here, is
\[
 \det O_K=5483119259214020183850397930008283625.
\]
The absence of a proper period for the vector would not, by itself, prove
invertibility: the decisive fact is that none of its Fourier modes
vanishes. The condition number is approximately \(8.15643\). If
\(K\) is normalized to unit norm, the extreme Gram eigenvalues are divided by
\(\|K\|^2=4251257\); the magnitude of its components must not
be confused with an intrinsic improvement over every reference of equal
norm.

\begin{corollary}[Polar separation of form and orientation]
\label{md:antmem:k:polar}
Let
\[
 G_K=O_K^*O_K,\qquad P_K=G_K^{1/2},\qquad
 F_K=O_KG_K^{-1/2}.
\]
Then \(P_K\) is positive definite, \(F_K\) is orthogonal, and
\(O_K=F_KP_K\). An isometric realization \(J:V\to\mathcal Y\)
preserves \(G_K\) and \(P_K\) and transports the isometric factor to
\(JF_K\).
\end{corollary}
\begin{proof}
Positivity of the Gram matrix allows its square root and inverse to be defined by the
spectral theorem. The identity
\(F_K^*F_K=G_K^{-1/2}G_KG_K^{-1/2}=I\) proves the claim.
Moreover, \((JO_K)^*(JO_K)=O_K^*O_K\).
\end{proof}

This polar decomposition is defined from the full Gram matrix: it cannot be obtained merely from
the list of eigenvalues or the norm of \(K\). The symbol
\(P_K\) in this corollary is not the rank-ten dimensional projector
of another incidence realization. This separation is a consequence of the
cyclic reference, not an identification of all its readout operators.

\subsection{Operator identification and relative holonomy}
\label{md:antmem:k:identificacion}

\begin{theorem}[Twelve general responses or one cyclic response]
\label{md:antmem:k:operadores}
For any linear operator \(H:V\to V\), its twelve responses
\(Y=(HK,HSK,\ldots,HS^{11}K)\) determine
\[
 H=YO_K^{-1}.
\]
If \(HS=SH\), a single complete vector response \(y=HK\)
determines
\begin{equation}
 H=O_yO_K^{-1},\qquad
 H=\sum_{j=0}^{11}h_jS^j,\qquad h=O_K^{-1}y.
 \label{md:antmem:k:filtro}
\end{equation}
With errors \(E\) in the response matrix or \(\delta y\) in the
vector, respectively,
\begin{align}
 \|\widehat H-H\|&\leq\frac{\|E\|}{\sqrt{589609}},
 \label{md:antmem:k:error-matriz}\\
 \|\widehat h-h\|&\leq\frac{\|\delta y\|}{\sqrt{589609}},\qquad
 \|\widehat H-H\|\leq\sqrt{\frac{12}{589609}}\,\|\delta y\|.
 \label{md:antmem:k:error-filtro}
\end{align}
\end{theorem}
\begin{proof}
The first equality is \(Y=HO_K\). If \(H\) commutes with \(S\),
then \(HS^jK=S^jy\), hence \(Y=O_y\). Commutation with
the cycle determines every column of \(H\) from a single one;
the twelve powers \(I,S,\ldots,S^{11}\) are independent and span
that operator space. Thus \(y=O_Kh\), giving the
second formula. The first two bounds use
\(\|O_K^{-1}\|=1/\sqrt{589609}\). The last uses
\(\|O_{\delta y}\|\leq\|O_{\delta y}\|_{\rm F}
=\sqrt{12}\|\delta y\|\).
\end{proof}

A response here means a twelve-component vector, not a scalar
measurement. For example, if \(H=2I+3S-S^5\), the response
\(y=HK\) returns exactly the coefficients \(2,3,-1\) in their
positions through \eqref{md:antmem:k:filtro}; the inversion procedure requires \(y,K,S\),
not the filter coefficients. By contrast, the nonzero operator whose first
row is \((543,-234,0,\ldots,0)\) and whose remaining rows vanish
annihilates \(K\). A single response does not identify an arbitrary operator.

The reference role admits an exact generalization. On a cycle of
\(n\) positions, \(H\mapsto Hg\) identifies every operator
commuting with the cycle if and only if
\(O_g=(g,Sg,\ldots,S^{n-1}g)\) is invertible. Sufficiency is the
preceding proof. If \(O_g\) is singular, a nonzero vector in its kernel
gives a polynomial of degree less than \(n\) that annihilates \(g\); the operator
is nonzero because the powers of the cycle are independent. The unit
impulse also has an orthonormal orbit and serves as a reference. The
specific consequence for \(K\) is that the same object already used in
closure and incidence satisfies that condition, without being selected
anew to optimize the test. No identification of nonlinear
dynamics or arbitrary hidden states follows; an application to
a linearization requires that it has been defined and justified.

\begin{corollary}[Recovery of relative holonomy from the complement]
\label{md:antmem:k:holonomia}
Let \(H_B\) be known and invertible, and let \(H_C\) be another transport of
the same type. The complement
\[
 D_\gamma=\frac{\sqrt8}{9}(H_C-H_B)
\]
determines
\[
 H_B^{-1}H_C=I+\frac9{\sqrt8}H_B^{-1}D_\gamma.
\]
If \(D_\gamma\) commutes with \(S\), in particular if both
transports do, \(y=D_\gamma K\) suffices to recover
\begin{equation}
 H_B^{-1}H_C
 =I+\frac9{\sqrt8}H_B^{-1}O_yO_K^{-1}.
 \label{md:antmem:k:holonomia-una-respuesta}
\end{equation}
The error in this reconstruction is bounded by
\[
 \frac9{\sqrt8}\|H_B^{-1}\|
 \sqrt{\frac{12}{589609}}\,\|\delta y\|.
\]
\end{corollary}
\begin{proof}
Solve for \(H_C\) in the definition of the complement and apply
\cref{md:antmem:k:operadores}. The bound follows by submultiplicativity.
\end{proof}

For \(H_B=S^3\) and \(H_C=S^4\), the rational datum
\(y/\sqrt8=(S^4-S^3)K/9\) makes it possible to reconstruct the relative holonomy
\(S\) without supplying \(H_C\) to the inversion procedure. If \(H_B\) is
unitary, the error factor is approximately \(0.01435510\). Without
commutation, the twelve general responses are retained, not a fictitious
symmetry assumption. Selection of a physical loop and its representation
still belong to its generating dynamics; this protocol does not
replace them.

\subsubsection{Consequence of the marked reference for two orientations}
\label{md:antmem:orientaciones}
The consequence used in the notebook retains the preceding
table as an antecedent. For a real operator \(H\) commuting with
\(S\), let \(y_+=HK\) and \(y_-=H^{\mathsf T}K\).
Their twelve Fourier intensities coincide, but \(y_+=y_-\) if and only
if \(H=H^{\mathsf T}\). Each complete vector response, together with
\(K\) and the orientation of \(S\), recovers its operator.
\begin{proof}
The operator \(H\) is circulant; its Fourier multipliers are
\(\lambda_m\), and those of \(H^{\mathsf T}\) are their conjugates.
Both responses therefore have intensities
\(|\lambda_m|^2|\widehat K_m|^2\).
If their vectors coincide, \((H-H^{\mathsf T})K=0\).
Commutation implies that \(H-H^{\mathsf T}\) also annihilates \(S^jK\)
for every \(j\). These vectors form a basis, so
\(H-H^{\mathsf T}=0\). Recovery is
\(H=O_{y_+}O_K^{-1}\), and analogously for \(H^{\mathsf T}\).
\end{proof}
In particular, the shifts \(S\) and \(S^{-1}\) give distinct responses
with equal power spectra. The exact autocorrelations
in the preceding table yield the identity already used in the notebook:
\begin{equation}
 \|SK-S^{-1}K\|^2
 =2\bigl(\|K\|^2-\langle K,S^2K\rangle\bigr)
 =2\,175\,744>0.
 \label{md:antmem:distancia-orientada}
\end{equation}
This number is a consequence of the autocorrelations, not an additional
starting datum. The notebook's rational check recovers from the two
responses the coefficients \(e_1\) and \(e_{11}\), respectively, by
inverting \(O_K\). A spectral intensity or scalar balance does not
replace the vector response.
This corollary concerns the twelve-component cyclic carrier:
it neither identifies \(S\) with the full enriched dynamics nor turns
\(S\leftrightarrow S^{-1}\) by itself into physical charge conjugation.
For general \(H\), transposition does not mean time reversal either;
for \(S\), it does coincide with \(S^{-1}\).

\subsection{Reversible coupling and nonstationary transport}
\label{md:antmem:acoplamiento}
\subsubsection{Reversible extension of the balance between two forms}

Let \(B,C:V\to V'\) be isomorphisms preserving a nondegenerate form \(b\) between two realizations:

\[
B^*b'B=C^*b'C=b.
\]

The form may be a positive metric or an alternating form. In real spaces, \(*\) denotes transposition with respect to the chosen coordinates; the equality expresses transport of the form between the specified realizations. In Hilbert spaces, the operators and their inverses are bounded.

With \(s=\sqrt8\) and \(Q_V=\frac13\left(\begin{smallmatrix}sI&-I\\I&sI\end{smallmatrix}\right)\), define

\begin{equation}
\boxed{
\mathcal U(B,C)=Q_{V'}^*\operatorname{diag}(B,C)Q_V
=\frac19\begin{pmatrix}8B+C&\sqrt8(C-B)\\\sqrt8(C-B)&B+8C\end{pmatrix}.}
\label{md:antmem:acop:eq:1}
\end{equation}

Write \(T=(8B+C)/9\), \(D=\sqrt8(C-B)/9\), and \(R=(B+8C)/9\). The full update is

\begin{equation}
\binom{x'}{m'}=\begin{pmatrix}T&D\\D&R\end{pmatrix}\binom{x}{m}.
\label{md:antmem:acop:eq:2}
\end{equation}

\begin{theorem}
The map \eqref{md:antmem:acop:eq:1} is reversible and preserves the form \(b\oplus b\):

\begin{equation}
\mathcal U(B,C)^*(b'\oplus b')\mathcal U(B,C)=b\oplus b,
\quad\mathcal U(B,C)^{-1}=\mathcal U(B^{-1},C^{-1}).
\label{md:antmem:acop:eq:3}
\end{equation}
\end{theorem}

\begin{proof}
The equality \(Q^*Q=I\) follows from \(8+1=9\). Its scalar blocks make \(Q\) preserve \(b\oplus b\) for any form \(b\). The diagonal map preserves the forms by hypothesis. Composition proves \eqref{md:antmem:acop:eq:3}, and inversion of the three factors gives the stated inverse.
\end{proof}

For incoming memory \(m=0\), the balance \(b=T^*b'T+D^*b'D\) is recovered. The second column determines the response to arbitrary incoming memory. Conservation of positive norms and conservation of oriented area are the respective specializations of \eqref{md:antmem:acop:eq:3}.

\paragraph{Collective realization in the nine copies}

In \(V^9\), let \(W_E=\operatorname{diag}(I,\ldots,I,E)\), with \(E\) in the ninth position. Let \(A(a,b)=(a/\sqrt8,\ldots,a/\sqrt8,b)\) and \(L=AQ\). The uniform embedding is \(Jx=(x,\ldots,x)/3\). Then \(L(x,0)=Jx\) and

\begin{equation}
W_E L=L\mathcal U(I,E).
\label{md:antmem:acop:eq:4}
\end{equation}

\begin{proof}
\(W_EA=A\operatorname{diag}(I,E)\); multiplication by \(Q\) and use of \(AQ=L\) give \eqref{md:antmem:acop:eq:4}. Orthogonality of the mean of the first eight positions to their seven differences also shows that the complement is fixed by \(W_E\).
\end{proof}

Let \(P_{\rm cyc}(v_0,\ldots,v_8)=(v_8,v_0,\ldots,v_7)\). The original cyclic step is \(S_E=P_{\rm cyc}W_E\). Thus \(S_EL=(P_{\rm cyc}L)U(I,E)\), with different input and output charts. The collective reduction preserves the corresponding pattern: \(S_E^9=I_9\otimes E\), whereas \(W_E^9=\operatorname{diag}(I_8,E^9)\). When a readout operator uses the nine individual labels, they are retained in their original representation.

\subsubsection{Exact composition and memory re-entry}

For \(B_1,C_1:V_0\to V_1\) and \(B_2,C_2:V_1\to V_2\),

\begin{equation}
\boxed{\mathcal U(B_2,C_2)\mathcal U(B_1,C_1)
=\mathcal U(B_2B_1,C_2C_1).}
\label{md:antmem:acop:eq:5}
\end{equation}

\begin{proof}
In \eqref{md:antmem:acop:eq:1}, the factors \(Q_{V_1}Q_{V_1}^*\) cancel, and the two diagonal blocks multiply in the indicated order. Induction extends the formula to any finite number of transports, preserving their order.
\end{proof}

The readout block contains the identity

\begin{equation}
\boxed{T_2T_1+D_2D_1=\frac{8B_2B_1+C_2C_1}{9}.}
\label{md:antmem:acop:eq:6}
\end{equation}

Starting from \((x,0)\), the first stage produces \((T_1x,D_1x)\). The second readout is \(T_2T_1x+D_2D_1x\). The second term is exactly the contribution of reintroduced memory. Storing \(D_1x\) outside the next coupling and continuing only with \(T_1x\) gives the separate-record procedure, whose readout is \(T_2T_1x\). Equality \eqref{md:antmem:acop:eq:6} specifies the difference between the two procedures.

The realization \(U\) preserves the coupling state and the ordered products \(B_N\cdots B_1\), \(C_N\cdots C_1\). Preservation of all prefixes also includes the HMT emission record, which retains each factor together with its order.

\paragraph{Covariant increment equation}

For the separate record \(x_{n+1}=T_nx_n\), \(m_n=D_nx_n\), one has

\begin{equation}
x_{n+1}-B_nx_n=m_n/\sqrt8.
\label{md:antmem:acop:eq:7}
\end{equation}

With \(G_0=I\) and \(G_{n+1}=B_nG_n\), telescoping gives

\begin{equation}
x_0=G_N^{-1}x_N-\frac1{\sqrt8}\sum_{n<N}G_{n+1}^{-1}m_n.
\label{md:antmem:acop:eq:8}
\end{equation}

Memory is the increment relative to the transport between forms \(B_n\), with the declared normalization. Equality \eqref{md:antmem:acop:eq:8} is a finite algebraic reconstruction; passing to an infinite series requires convergence. The following theorem constructs the stable inverse through the limit of positive operators.

\subsubsection{Memory of a nonstationary sequence}

In a Hilbert-space realization \(H\), let \(C_n\) be unitary, with their order of composition declared. Write

\[
T_n=(8I+C_n)/9,\quad D_n=\sqrt8(C_n-I)/9,
\quad P_0=I,\quad P_{n+1}=T_nP_n,\quad A_N=P_N^*P_N.
\]

The index \(N\) counts compressions with separate storage, unlike full transport with re-entry. For variable metrics with isometric isomorphisms \(B_n\), conjugation by \(G_n\) takes the transports to a fixed Hilbert space and yields this same case, with \(\widetilde C_n=G_{n+1}^{-1}C_nG_n\). The statement uses the fibre metrics and the declared isometries between them.

\begin{theorem}[Positive terminal and nonstationary recovery]
\label{md:antmem:terminal-no-estacionario}
There is a positive operator \(0\leq A_\infty\leq I\), the strong limit of \(A_N\), and

\begin{equation}
\boxed{I=A_\infty+\sum_{n\ge0}P_n^*D_n^*D_nP_n.}
\label{md:antmem:acop:eq:9}
\end{equation}

The map

\begin{equation}
\boxed{\mathcal Vv=(A_\infty^{1/2}v,(D_nP_nv)_{n\ge0})}
\label{md:antmem:acop:eq:10}
\end{equation}

is isometric and has a stable left inverse

\begin{equation}
\mathcal V^*(a,(m_n))=A_\infty^{1/2}a+
\sum_{n\ge0}P_n^*D_n^*m_n,\qquad\|\mathcal V^*\|=1.
\label{md:antmem:acop:eq:11}
\end{equation}

\end{theorem}
\begin{proof}
The local identity \(T_n^*T_n+D_n^*D_n=I\) gives
\(A_N-A_{N+1}=P_N^*D_N^*D_NP_N\geq0\). The quadratic forms of
the decreasing positive sequence converge to that of a positive operator
\(A_\infty\). For \(B_N=A_N-A_\infty\), the inequality
\(0\leq B_N^2\leq B_N\) turns this convergence into strong convergence.
Finite telescoping and passage to the limit prove \eqref{md:antmem:acop:eq:9}. Evaluation on \(v\)
proves \eqref{md:antmem:acop:eq:10}. Its adjoint is \eqref{md:antmem:acop:eq:11}; the series converges in norm because
the truncations of a square-summable record converge and the adjoint is bounded.
\end{proof}

Reconstruction with the limiting terminal and \(N\) records has exact
error \((A_N-A_\infty)v\). The vector \(A_\infty^{1/2}v\) is the
canonical positive realization of the norm terminal. Its existence follows
from the Gram limit, independently of vector convergence of \(P_Nv\).

\subsubsection{Three different properties: exhaustion, distinguishability, and stability}

Let \(Mv=(D_nP_nv)_n\). By \eqref{md:antmem:acop:eq:9}, \(M^*M=I-A_\infty\). Hence:

\begin{enumerate}
\item The entire norm of \(v\) passes to the record exactly when \(A_\infty v=0\).
\item The record distinguishes every state exactly when \(\ker(I-A_\infty)=\{0\}\).
\item Recovery from memory alone has a bounded inverse exactly when
\(I-A_\infty\geq\varepsilon I\) for some \(\varepsilon>0\).
In that case a left inverse is \((I-A_\infty)^{-1}M^*\).
\end{enumerate}

Moreover,

\begin{equation}
\boxed{\ker M=\bigcap_{n\ge0}\operatorname{Fix}(C_n).}
\label{md:antmem:acop:eq:12}
\end{equation}

\begin{proof}
If \(Mv=0\), then \(D_0v=0\), whence \(C_0v=v\)
and \(P_1v=v\). Inductively, \(D_nP_nv=0\) gives \(C_nv=v\) and
\(P_{n+1}v=v\). The converse is immediate. The three equivalences
follow from the Gram operator \(M^*M\) and the characterization of
operators bounded below.
\end{proof}

\paragraph{Exact example separating the three properties}

In the algebraic example \(C_0=-I\) and \(C_n=I\) for \(n\geq1\),

\begin{equation}
A_\infty=\frac{49}{81}I,\qquad M^*M=\frac{32}{81}I,
\qquad m_0=-\frac{2\sqrt8}{9}v,
\qquad v=-\frac9{2\sqrt8}m_0.
\label{md:antmem:acop:eq:13}
\end{equation}

Memory contains \(32/81\) of the squared norm and permits recovery of the entire vector. The terminal retains \(49/81\) and is a positive operator distinct from a projector. Recoverability depends on the kernel and conditioning of the readout operator. Readout and memory are correlated components of an isometry.

This example separates the theorem's three logical consequences. An
application to a physical trajectory also preserves the effective selection and
provenance of its \(C_n\) through TPK.

The letter \(\eta\) in the following Gaussian realization denotes
the deformation of the ellipse, namely the parameter \(\zeta\) of the
transport developments; it is not the action coefficient
\(\eta_{\rm ret}\). A positive section \(\hbar_a\) and
the same unit for both coordinates of the angular ratio are maintained.
The retained semiaxes are
\[
 a_S=\sqrt{2\hbar_a}\,e^{\eta/2},\qquad
 b_S=\sqrt{2\hbar_a}\,e^{-\eta/2},\qquad
 \eta=\operatorname{artanh}(C^*/A).
\]

\subsection{Canonical covariance and reciprocity of the action ellipse}
\label{md:antmem:covarianza-accion}

The preceding construction determines two complementary invariants:
the product of the semiaxes fixes the action area, and their ratio retains
the anisotropy of the angular pair. We now specify the canonical realization
in which those semiaxes are twice the standard deviations. The starting
structure is the ellipse already obtained from the angular and action
outputs of APP--TRIT--TPK; the operator realization is established after
that construction.

In this subsection, write \(\eta=\eta_{\rm el}
=\operatorname{artanh}(C^*/A)\), with \(|C^*/A|<1\) and
\(\hbar_a>0\). The angular ratio uses the same unit for both
components. In balanced coordinates \((Q,P)\), both having the dimensions
of the square root of action, define
\begin{equation}
 V_\eta:=\frac14
 \begin{pmatrix}a_S^2&0\\0&b_S^2\end{pmatrix}
 =\frac{\hbar_a}{2}
 \begin{pmatrix}e^\eta&0\\0&e^{-\eta}\end{pmatrix}.
 \label{md:antmem:covarianza-geometrica}
\end{equation}
This positive matrix has entries having the dimensions of action. Its interpretation
as a quantum covariance is specified by the following realization.

\begin{proposition}[Gaussian realization of the action covariance]
\label{md:antmem:covarianza-gaussiana}
Consider the canonical representation on \(L^2(\mathbb R,dQ)\), with
\(\widehat Qf(Q)=Qf(Q)\) and
\(\widehat Pf(Q)=-i\hbar_a f'(Q)\), initially on the Schwartz
space \(\mathcal S(\mathbb R)\). The normalized state
\begin{equation}
 \psi_\eta(Q)=\bigl(\pi\hbar_a e^\eta\bigr)^{-1/4}
 \exp\!\left(-\frac{Q^2}{2\hbar_a e^\eta}\right)
 \label{md:antmem:gaussiana-accion}
\end{equation}
has zero means and covariance matrix \(V_\eta\). In particular,
\begin{equation}
 (\Delta Q)^2=\frac{\hbar_a e^\eta}{2},\qquad
 (\Delta P)^2=\frac{\hbar_a e^{-\eta}}{2},\qquad
 \operatorname{Cov}(Q,P)=0,\qquad
 \det V_\eta=\frac{\hbar_a^2}{4}.
 \label{md:antmem:varianzas-accion}
\end{equation}
This realization saturates the Robertson--Schrödinger inequality.
\end{proposition}

\begin{proof}
The preceding operators are essentially self-adjoint on
\(\mathcal S(\mathbb R)\), which is a common invariant core;
on this space,
\([\widehat Q,\widehat P]=i\hbar_a I\). For
\(\widehat Y=(\widehat Q,\widehat P)\), the symmetric covariance is
\[
 V_{jk}=\frac12\left\langle
 (\widehat Y_j-\langle\widehat Y_j\rangle)
 (\widehat Y_k-\langle\widehat Y_k\rangle)
 +(\widehat Y_k-\langle\widehat Y_k\rangle)
 (\widehat Y_j-\langle\widehat Y_j\rangle)
 \right\rangle.
\]
With \(s=\hbar_a e^\eta>0\), let
\(I_s=\int_{\mathbb R}e^{-Q^2/s}\,dQ\). The product of two
equal integrals, evaluated in polar coordinates, gives
\(I_s^2=2\pi\int_0^\infty r e^{-r^2/s}\,dr=\pi s\).
Moreover, the integral of
\(\frac{d}{dQ}(Qe^{-Q^2/s})\) is zero because of decay at both
ends; hence
\(\int_{\mathbb R}Q^2e^{-Q^2/s}\,dQ=sI_s/2\).
These identities and parity prove normalization, the zero mean, and
\(\langle\widehat Q^2\rangle=s/2\). Moreover,
\(\widehat P\psi_\eta=i e^{-\eta}Q\psi_\eta\).
Thus \(\langle\widehat P\rangle=0\),
\(\|\widehat P\psi_\eta\|^2=\hbar_a e^{-\eta}/2\), and
\(\operatorname{Re}\langle\widehat Q\psi_\eta,
\widehat P\psi_\eta\rangle=0\).

For any normalized state in that domain, Cauchy--Schwarz applied
to \((\widehat Q-\langle\widehat Q\rangle)\psi\) and
\((\widehat P-\langle\widehat P\rangle)\psi\) gives
\[
 (\Delta Q)^2(\Delta P)^2
 \geq \operatorname{Cov}(Q,P)^2
       +\frac14\left|\langle[\widehat Q,\widehat P]\rangle\right|^2
 =\operatorname{Cov}(Q,P)^2+\frac{\hbar_a^2}{4}.
\]
The computed variances attain equality. Equivalently, the
operator
\[
 a_\eta=\frac{e^{-\eta/2}\widehat Q
                  +i e^{\eta/2}\widehat P}{\sqrt{2\hbar_a}}
\]
satisfies \([a_\eta,a_\eta^\dagger]=I\) and
\(a_\eta\psi_\eta=0\), by direct substitution.
\end{proof}

\subsection{Covariance, normal form, and Gaussian spectrum}
\label{md:antmem:gauss:gaussianas-espectro}

The positive action section \(\hbar_a\) and the preceding canonical
realization determine the normalization used here. We prove the
correspondence between covariance, spectrum, purity, and entropy for a
finite number of modes; Fock space retains all its occupation numbers.
Symplectic reduction and the Gaussian formulae are established
results reproduced here to make their application to memory
transport self-contained.

\subsubsection{Canonical coordinates and characteristic function}
On \(L^2(\mathbb R^d)\), let \(\widehat q_j=\widehat Q_j/\sqrt{\hbar_a}\),
\(\widehat p_j=\widehat P_j/\sqrt{\hbar_a}\), and
\(\widehat z=(\widehat q_1,\widehat p_1,\ldots,\widehat q_d,\widehat p_d)\).
On Schwartz space,
\[
 [\widehat z_j,\widehat z_k]=iJ_{jk}I,\qquad
 J=\bigoplus_{j=1}^d\begin{pmatrix}0&1\\-1&0\end{pmatrix}.
\]
This convention fixes the sign of the canonical form. The Weyl
operators \(\mathcal W(\xi)=\exp(i\xi^{\mathsf T}\widehat z)\) are defined
through the self-adjoint closures of the indicated linear combinations.
For a centred state \(\rho\) with finite second moments, the
normalized covariance and its characteristic function are
\[
 W_{jk}=\operatorname{Tr}\rho(\widehat z_j\widehat z_k+\widehat z_k\widehat z_j),
 \qquad \chi_\rho(\xi)=\operatorname{Tr}(\rho\mathcal W(\xi)),\qquad
 V=\frac{\hbar_a}{2}W.
\]
The state is centred Gaussian when
\begin{equation}
 \chi_\rho(\xi)=\exp(-\xi^{\mathsf T}W\xi/4).
 \label{md:antmem:gauss:eq:caracteristica-gaussiana}
\end{equation}
The expectation of \((c^{\mathsf T}\widehat z)^\dagger
(c^{\mathsf T}\widehat z)\) is \(c^\dagger(W+iJ)c/2\); thus
\(W+iJ\geq0\). Moreover, \(W>0\): if \(v\) is real and
\(v^{\mathsf T}Wv=0\), positivity implies \((W+iJ)v=0\),
hence \(Jv=0\) and \(v=0\).

\begin{lemma}[Uniqueness of the characteristic function]
\label{md:antmem:gauss:lem:unicidad-weyl}
Two trace-class operators with the same characteristic function are equal.
\end{lemma}
\begin{proof}
Regrouping positions and momenta as \(\xi=(u,v)\),
\[
 (\mathcal W(u,v)f)(x)=e^{iu\cdot(x+v/2)}f(x+v).
\]
For a finite-rank operator with Schwartz vectors and kernel \(K_A\),
\[
 \operatorname{Tr}(A\mathcal W(u,v))=\int K_A(x+v/2,x-v/2)e^{iu\cdot x}\,dx.
\]
Plancherel's theorem and the change of variables with absolute Jacobian one give
\[
 \|A\|_{\rm HS}^2=(2\pi)^{-d}\int_{\mathbb R^{2d}}
 |\operatorname{Tr}(A\mathcal W(\xi))|^2\,d\xi.
\]
These operators are dense in the trace class. Convergence
in that norm implies Hilbert--Schmidt convergence and uniform
convergence of the characteristic functions, since their difference is
bounded by \(\|A-B\|_1\). The equality extends in \(L^2\).
Applying it to the difference of the operators proves the claim.
\end{proof}

\subsubsection{Symplectic reduction of the positive form}
\begin{theorem}[Williamson normal form]
\label{md:antmem:gauss:thm:williamson}
For a real symmetric matrix \(W>0\) of order \(2d\), there exist
\(S\in\operatorname{Sp}(2d,\mathbb R)\) and \(\nu_j>0\) such that
\[
 W=SDS^{\mathsf T},\qquad D=\bigoplus_j\nu_jI_2.
\]
The eigenvalues of \(JW\) are \(\pm i\nu_j\). The condition
\(W+iJ\geq0\) is equivalent to \(\nu_j\geq1\) for every \(j\).
\end{theorem}
\begin{proof}
The matrix \(B=W^{1/2}JW^{1/2}\) is skew-symmetric and invertible.
If \(-B^2u=\nu^2u\) and \(\|u\|=1\), the vector
\(v=-Bu/\nu\) has unit norm and is orthogonal to \(u\); it satisfies
\(Bu=-\nu v\), \(Bv=\nu u\). The complement of this pair is
invariant. Iteration produces an orthogonal matrix \(O\) with
\(O^{\mathsf T}BO=JD\). Defining \(S=W^{1/2}OD^{-1/2}\),
\[
 S^{\mathsf T}JS=J,\qquad SDS^{\mathsf T}=W.
\]
The similarity \(W^{1/2}(JW)W^{-1/2}=B\) determines the values
\(\nu_j\), with multiplicity. Congruence by \(S^{-1}\)
transforms \(W+iJ\) into \(D+iJ\); each block has eigenvalues
\(\nu_j+1\) and \(\nu_j-1\).
\end{proof}

\begin{lemma}[Unitary implementation in finitely many modes]
\label{md:antmem:gauss:lem:implementacion-simplectica}
Every symplectic matrix \(S\) admits a unitary \(U_S\) satisfying
\(U_S^\dagger\mathcal W(\xi)U_S=\mathcal W(S^{\mathsf T}\xi)\).
The conjugation \(\rho\mapsto U_S\rho U_S^\dagger\) transports
covariance by \(W\mapsto SWS^{\mathsf T}\).
\end{lemma}
\begin{proof}
In the polar decomposition \(S=OP\), applying the identity
\(J^{-1}(S^{\mathsf T}S)J=(S^{\mathsf T}S)^{-1}\) to the positive square root
through spectral calculus proves \(PJP=J\).
Then \(O\) is orthogonal symplectic. The eigenspaces of \(P\)
are paired by \(J\) with eigenvalues \(\lambda,\lambda^{-1}\);
the eigenspace for eigenvalue one is \(J\)-invariant. An
orthogonal symplectic matrix \(K\) therefore reduces \(P\) to
\(\bigoplus_j\operatorname{diag}(e^{r_j},e^{-r_j})\).

Dilation is implemented by
\[
 (\mathcal D_rf)(x)=e^{-\frac12\sum_jr_j}
 f(e^{-r_1}x_1,\ldots,e^{-r_d}x_d).
\]
A change of variables proves unitarity and the conjugation of
\((\widehat q_j,\widehat p_j)\) to
\((e^{r_j}\widehat q_j,e^{-r_j}\widehat p_j)\).
An orthogonal symplectic transformation commutes with \(J\) and
corresponds to a matrix \(u\in U(d)\) acting on
\(a_j=(\widehat q_j+i\widehat p_j)/\sqrt2\). Its second
quantization \(\Gamma(u)=\bigoplus_{n\geq0}u^{\otimes_s n}\)
is unitary on \(\mathcal F_s(\mathbb C^d)\), and the identity on
finite symmetric products gives the transformation of the \(a_j\).

The Hermite basis identifies this space with \(L^2(\mathbb R^d)\).
To justify completeness, if \(f\) is orthogonal to every
polynomial multiplied by \(e^{-|x|^2/2}\), the entire function
\(F(z)=\int f(x)e^{-|x|^2/2}e^{z\cdot x}\,dx\) has all its
derivatives at zero equal to zero. Thus \(F=0\); restricting to
\(z=i\xi\), Fourier uniqueness gives \(f=0\). The
normalizations and orthogonality follow from creation and annihilation.
Composing the preceding implementations produces \(U_S\).
The Weyl equality implies
\(\chi_{U_S\rho U_S^\dagger}(\xi)=\chi_\rho(S^{\mathsf T}\xi)\),
which proves the covariance law.
\end{proof}

\subsubsection{Computing the spectrum of the isotropic density}
In one mode, the normalized coherent states are
\[
 |z\rangle=e^{-|z|^2/2}\sum_{n\geq0}\frac{z^n}{\sqrt{n!}}|n\rangle.
\]
Their position wavefunction,
\(\pi^{-1/4}\exp[-(x-q_0)^2/2+ip_0(x-q_0/2)]\), where
\(q_0=\sqrt2\Re z\), \(p_0=\sqrt2\Im z\), gives by integration
\[
 \langle z|\mathcal W(u,v)|z\rangle
 =e^{-(u^2+v^2)/4}e^{i(uq_0+vp_0)}.
\]
For \(s>0\), the integral
\[
 \rho^{(s)}=\frac1{\pi s}\int_{\mathbb C} e^{-|z|^2/s}
 |z\rangle\langle z|\,d^2z
\]
converges in trace norm, is positive, and has trace one. Angular
integration annihilates the off-diagonal elements in the occupation basis;
radial integration gives
\[
 \langle n|\rho^{(s)}|n\rangle=\frac{s^n}{(1+s)^{n+1}},\qquad
 \chi_{\rho^{(s)}}(u,v)=e^{-(2s+1)(u^2+v^2)/4}.
\]
With \(\nu=2s+1\), the unique Gaussian state with covariance
\(\nu I_2\) is, by Lemma~\ref{md:antmem:gauss:lem:unicidad-weyl},
\begin{equation}
 \rho_\nu=(1-t)\sum_{n\geq0}t^n|n\rangle\langle n|,
 \qquad t=\frac{\nu-1}{\nu+1}.
 \label{md:antmem:gauss:eq:estado-geometrico}
\end{equation}
The case \(\nu=1\) is the vacuum projector.

\begin{theorem}[Spectrum, purity, and entropy]
\label{md:antmem:gauss:thm:espectro-gaussiano}
A centred Gaussian state with covariance \(W\) and symplectic
eigenvalues \(\nu_1,\ldots,\nu_d\) is unitarily equivalent
to \(\bigotimes_j\rho_{\nu_j}\). Its eigenvalues are
\[
 \lambda_{\mathbf n}=\prod_j(1-t_j)t_j^{n_j},\qquad
 t_j=\frac{\nu_j-1}{\nu_j+1},\quad \mathbf n\in\mathbb N^d.
\]
With \(0\log0=0\),
\begin{align}
 \operatorname{Tr}\rho^2&=\prod_j\nu_j^{-1}=(\det W)^{-1/2},
 \label{md:antmem:gauss:eq:pureza-general}\\
 -\operatorname{Tr}(\rho\log\rho)&=\sum_j
 \left[\frac{\nu_j+1}{2}\log\frac{\nu_j+1}{2}
 -\frac{\nu_j-1}{2}\log\frac{\nu_j-1}{2}\right].
 \label{md:antmem:gauss:eq:entropia-general}
\end{align}
\end{theorem}
\begin{proof}
Williamson's theorem gives \(W=SDS^{\mathsf T}\). The product of the isotropic
densities has characteristic function \(e^{-\xi^{\mathsf T}D\xi/4}\).
Unitary implementation and Weyl uniqueness prove the
equivalence and the spectral formula. In one mode,
\[
 \sum_{n\geq0}(1-t)^2t^{2n}=\frac{1-t}{1+t}=\nu^{-1},
\]
and
\[
 -\sum_{n\geq0}(1-t)t^n\log((1-t)t^n)
 =-\log(1-t)-\frac{t}{1-t}\log t.
\]
Substitution of \(t=(\nu-1)/(\nu+1)\) proves the formulae; the
limit at \(t=0\) is zero. In the finite product, Tonelli's theorem justifies
summing the nonnegative entropy terms, and
\(\det W=\prod_j\nu_j^2\) gives the determinantal expression.
\end{proof}

If the state is the reduction of a pure bipartite vector, its
eigenvalues are the squares of the Schmidt coefficients.
The singular-value decomposition of the Hilbert--Schmidt operator defined
by that vector proves equality of the nonzero spectra of both
reductions. The entropy computed is then entanglement entropy.

\subsubsection{Preparation and transport of the two covariances}
The subsequent application uses two identical pure Gaussian inputs.
In general, for \(M\in\operatorname{Sp}(2d,\mathbb R)\), their covariance is
\(V_0=(\hbar_a/2)MM^{\mathsf T}\), and the joint input has
covariance \(V_0\oplus V_0\). The physical transport is obtained by
conjugating \(U_H\) by \(M\oplus M\). In the balanced chart,
the normalized input covariance is \(I_{4d}\); with
\[
 T=\frac{8I+H}{9},\qquad D=\frac{\sqrt8(H-I)}9,
\]
the first reduction has
\[
 W=TT^{\mathsf T}+DD^{\mathsf T}=\frac{8I+HH^{\mathsf T}}9.
\]
Its characteristic function is obtained by setting the arguments of the
second group of modes to zero. It therefore remains Gaussian, and the preceding
theorems compute its entire spectrum without truncating the occupation record.

\subsection{Prolongation of memory transport: quantum realization}

The quantum realization uses the reversible coupling of Section~\ref{md:antmem:acoplamiento}. We first consider a real canonical plane with alternating matrix \(\Omega=\left(\begin{smallmatrix}0&1\\-1&0\end{smallmatrix}\right)\). The nonadic partition determines

\begin{equation}
Q=\frac13\begin{pmatrix}\sqrt8 I&-I\\I&\sqrt8 I\end{pmatrix},\qquad
U_H=Q^{\mathsf T}\operatorname{diag}(I,H)Q
=\frac19\begin{pmatrix}8I+H&\sqrt8(H-I)\\
\sqrt8(H-I)&I+8H\end{pmatrix}.
\label{md:antmem:cuantica:eq:5.1}
\end{equation}

The operator \(H\in\operatorname{Sp}(2,\mathbb R)\) transports that plane. The matrix \(Q\) preserves the Euclidean inner product and the form \(\Omega\oplus\Omega\); hence \(U_H\) is symplectic. The first component defines the observed system and the second the memory record; the full transformation acts on both.

The ordered loop whose identities are reproduced in
\ref{md:antmem:lazo-ep} is

\begin{equation}
C=\begin{pmatrix}0&-1\\1&-1\end{pmatrix},\quad
P=\begin{pmatrix}1&1\\0&1\end{pmatrix},\quad
H_n=C^{-1}P^{-n}CP^n
=\begin{pmatrix}1+n&n^2\\n&n^2-n+1\end{pmatrix}.
\label{md:antmem:cuantica:eq:5.2}
\end{equation}

For every \(n\in\mathbb Z\), \(\det H_n=1\). In particular, for \(n=1\),

\begin{equation}
H_1=\begin{pmatrix}2&1\\1&1\end{pmatrix}
=F_1^2,\quad F_1=\begin{pmatrix}1&1\\1&0\end{pmatrix}
=P F_{\rm av}P^{-1}.
\label{md:antmem:cuantica:eq:5.3}
\end{equation}

Its spectrum is \(\{\varphi^2,\varphi^{-2}\}\), by subsequent recognition of the golden mode already generated. The integer \(n\) specifies the oriented number of iterations in the chosen loop. The collection of transports is algebraically determined; its realization as a physical protocol also retains the corresponding preparation and evolution data.

\subsubsection{Algebraic identities of the loop and subsequent recognition}
\label{md:antmem:lazo-ep}
To preserve the proof of the identities used, their matrix
composition is reproduced here. In the same chart,
\[
 N_0=\begin{pmatrix}0&1\\0&0\end{pmatrix}.
\]
Let
\[
 P=I+N_0=\begin{pmatrix}1&1\\0&1\end{pmatrix}.
\]
Since \(N_0^2=0\), one has \(P^n=I+nN_0\) for every
\(n\in\mathbb Z\).
The ordered loop of four transports gives

\begin{equation}
H_n=C^{-1}P^{-n}CP^n
=\begin{pmatrix}1+n&n^2\\n&n^2-n+1\end{pmatrix}.
\label{md:antmem:matriz:eq:28}
\end{equation}

For \(n\ne0\), its normalized difference is

\[
F_n=\frac{H_n-I}{n}
=\begin{pmatrix}1&n\\1&n-1\end{pmatrix}.
\]

Multiplying these matrices yields, without introducing spectral values,

\begin{equation}
\boxed{F_n^2=H_n,\qquad F_n^2-nF_n-I=0.}
\label{md:antmem:matriz:eq:29}
\end{equation}

With \(H_B=I\), the memory is
\(D_{\gamma,n}=(\sqrt8/9)nF_n\).
Consequently,

\begin{equation}
\boxed{H_n=\left(\frac9{n\sqrt8}D_{\gamma,n}\right)^2.}
\label{md:antmem:matriz:eq:30}
\end{equation}

The matrix memory response reconstructs the hyperbolic
self-scale. For \(n>0\), subsequent spectral recognition
identifies

\begin{equation}
\rho_{{\rm ep},n}=\frac{n+\sqrt{n^2+4}}2,\quad
\operatorname{Spec}F_n=\{\rho_{{\rm ep},n},-\rho_{{\rm ep},n}^{-1}\},\quad
\operatorname{Spec}H_n=\{\rho_{{\rm ep},n}^2,\rho_{{\rm ep},n}^{-2}\}.
\label{md:antmem:matriz:eq:31}
\end{equation}

The case \(n=1\) produces the golden polynomial.
The relation to the operator
\(F_{\rm av}=\begin{pmatrix}0&1\\1&1\end{pmatrix}\)
generated by the calendar preserves the change
of chart:

\begin{equation}
F_1=PF_{\rm av}P^{-1},\qquad
H_1=PF_{\rm av}^{\,2}P^{-1}.
\label{md:antmem:matriz:eq:32}
\end{equation}

The transport \(P\), of determinant one, intertwines the two
matrices and preserves \(\Omega\).

\subsubsection{Initial state, transport, and readout}

Proposition~\ref{md:antmem:covarianza-gaussiana} constructs the Gaussian realization of the action ellipse. On the positive section \(\hbar_a>0\), with \(A>0\) and \(|C^*/A|<1\), its covariance is

\begin{equation}
V_\eta=\frac{\hbar_a}{2}M_\eta M_\eta^{\mathsf T},\qquad
M_\eta=\operatorname{diag}(e^{\eta/2},e^{-\eta/2}),\qquad
\eta=\operatorname{artanh}(C^*/A).
\label{md:antmem:cuantica:eq:5.4}
\end{equation}

In the Schrödinger representation on \(L^2(\mathbb R,dQ)\), with \(\widehat Q\) multiplication by \(Q\) and \(\widehat P=-i\hbar_a\,d/dQ\), the relation \([\widehat Q,\widehat P]=i\hbar_a I\) holds on Schwartz space. The normalized pure state is

\[
\psi_\eta(Q)=(\pi\hbar_a e^\eta)^{-1/4}
\exp[-Q^2/(2\hbar_a e^\eta)].
\]

The preparation is the product \(\psi_\eta\otimes\psi_\eta\). The transport in \eqref{md:antmem:cuantica:eq:5.1} is expressed in the same covariance chart by

\begin{equation}
\widetilde U=(M_\eta\oplus M_\eta)U_H
(M_\eta^{-1}\oplus M_\eta^{-1}).
\label{md:antmem:cuantica:eq:5.5}
\end{equation}

Lemma~\ref{md:antmem:gauss:lem:implementacion-simplectica} provides a unitary implementation of this symplectic matrix in two modes. Implementations differing by a global phase produce the same conjugation of the state. Covariance transforms by \(V\mapsto\widetilde U V\widetilde U^{\mathsf T}\); the joint state remains pure, and its partial trace over the second mode determines the visible readout, in accordance with the construction of Section~\ref{md:antmem:gauss:gaussianas-espectro}.

\begin{theorem}[Mixedness of the reduced state from relative deformation]

The covariance of the visible readout is

\begin{equation}
\boxed{V_{\rm vis}=
\frac{\hbar_a}{2}M_\eta
\frac{8I+HH^{\mathsf T}}9
M_\eta^{\mathsf T}.}
\label{md:antmem:cuantica:eq:5.6}
\end{equation}

Its normalized symplectic eigenvalue is

\begin{equation}
\boxed{\nu(H):=\frac{2\sqrt{\det V_{\rm vis}}}{\hbar_a}
=\sqrt{1+\frac8{81}\left[\operatorname{tr}(HH^{\mathsf T})-2\right]}.}
\label{md:antmem:cuantica:eq:5.7}
\end{equation}
\end{theorem}

\begin{proof}
Let \(T=(8I+H)/9\) and \(D=\sqrt8(H-I)/9\). The two initial modes are independent and have the same covariance. In the balanced chart of the initial state, the normalized visible covariance is \(TT^{\mathsf T}+DD^{\mathsf T}\). Expansion gives

\[
\begin{aligned}
TT^{\mathsf T}+DD^{\mathsf T}
&=\frac{64I+8(H+H^{\mathsf T})+HH^{\mathsf T}
+8[HH^{\mathsf T}-H-H^{\mathsf T}+I]}{81}\\
&=\frac{8I+HH^{\mathsf T}}9.
\end{aligned}
\]

This proves \eqref{md:antmem:cuantica:eq:5.6}. Since \(\det M_\eta=1\) and \(\det(HH^{\mathsf T})=1\),
\(\det(8I+HH^{\mathsf T})=65+8\operatorname{tr}(HH^{\mathsf T})\).
Division by \(81\) proves \eqref{md:antmem:cuantica:eq:5.7}. The two positive eigenvalues of \(HH^{\mathsf T}\) are reciprocal; their sum is at least two. Consequently, \(\nu\geq1\), with equality exactly when \(HH^{\mathsf T}=I\).
\end{proof}

The scale \(\hbar_a\) and deformation \(\eta\) are retained until the determinant is taken. Their cancellation is covariant because state and operator are transported jointly. The expression \(HH^{\mathsf T}\) uses the Euclidean inner product of the balanced chart of the initial state; a change of chart also transports that metric inner product.

An orthogonal symplectic transport \(H\) may have \(D\neq0\) while preserving the isotropic Gaussian product, with \(\nu=1\). The operator complement and state entanglement therefore have different criteria. In this preparation, nonisometric relative deformation characterizes exactly the case \(\nu>1\).

\subsubsection{Exact evaluation of the golden loop}

For \eqref{md:antmem:cuantica:eq:5.2},

\begin{equation}
\operatorname{tr}(H_nH_n^{\mathsf T})-2
=n^2(2n^2-2n+5),\qquad
\nu_n^2=1+\frac8{81}n^2(2n^2-2n+5).
\label{md:antmem:cuantica:eq:5.8}
\end{equation}

The identity follows by summing the squares of the four entries. For \(n=1\), the sum is seven, and

\begin{equation}
\boxed{\nu_1=\frac{11}{9},\quad
\operatorname{Tr}\rho_{\rm vis}^2=\frac9{11},\quad
\bar n_{\rm W}=\frac19,\quad
\frac{S_{\rm vis}}{k_B}=\frac{10}{9}\log10-\log9.}
\label{md:antmem:cuantica:eq:5.9}
\end{equation}

Theorem~\ref{md:antmem:gauss:thm:espectro-gaussiano} proves that the centred one-mode Gaussian state with symplectic eigenvalue \(\nu\) is unitarily equivalent to the geometric density in \eqref{md:antmem:gauss:eq:estado-geometrico}, whose mean occupation is \(\bar n_{\rm W}=(\nu-1)/2\). For \(\nu=11/9\), its eigenvalues are

\begin{equation}
\lambda_j=\frac9{10}\,10^{-j},\qquad j=0,1,2,\ldots.
\label{md:antmem:cuantica:eq:5.10}
\end{equation}

The series sums to one. The sum of squares is
\((9/10)^2/(1-10^{-2})=9/11\). Finally,

\[
-\sum_j\lambda_j\log\lambda_j
=-\log(9/10)+\log10\sum_jj\lambda_j
=\frac{10}{9}\log10-\log9.
\]

Gaussian reduction and its spectrum are proved in Section~\ref{md:antmem:gauss:gaussianas-espectro}; the preceding two series evaluate exactly the purity and entropy of this realization. Since the joint state is pure, \(S_{\rm vis}/k_B\) is also its dimensionless entanglement entropy. The global transformation is unitary, whereas restriction to the first mode has positive entropy.

The result concerns the specified preparation and transport. Identifying it with a vacuum measurement requires the experimental realization of that protocol. The entropy of a reduction and transferred heat remain distinct quantities.
